%============================================================
%  Bd. 41 - Abelsche Gruppen
%============================================================
\documentclass[main.tex]{subfiles}
\ifSubfilesClassLoaded{\BandSetupStandalone{B41}}{}
\title{Bd. 41 - Abelsche Gruppen}
\author{Martin Kunze}
\date{}
\setcounter{file}{41}
\hypersetup{pdftitle={Bd. 41 - Abelsche Gruppen},pdfauthor={Martin Kunze}}
\begin{document}
\providecommand{\BandXXIXProofRefStack}[1]{%
  \ensuremath{\begin{gathered}[t]#1\end{gathered}}}
\title{Bd. 41 - Abelsche Gruppen}
\author{Martin Kunze}
\date{}
\hypersetup{pageanchor=false}
\maketitle
\hypersetup{pageanchor=true}
\setcounter{tocdepth}{3}
\tableofcontents
\setcounter{file}{41}
\renewcommand{\FormulaBandID}{41}

\chapter{Einleitung}

Dieser Band spezialisiert die Gruppen aus Band~40 durch das
Kommutativitätsgesetz. Neben den zugehörigen Morphismen behandelt er
die formalen Differenzen eines kommutativen kürzbaren Monoids.
Die Quotientenkonstruktion liefert eine abelsche Gruppe und eine
injektive Monoidabbildung des Ausgangsmonoids. Die bereits in Band~17
konstruierten ganzen Zahlen bilden das grundlegende additive Beispiel.

\chapter{Abelsche Gruppen}

\FormulaAxiomDeltaK[Kommutativität einer Gruppenoperation]{%
  \GroupAlg{G}{\circ}{e}\dsep x,y\in G
  \vdash x\circ y=y\circ x%
}{AbelianGroupCommutativityLaw}{%
  \DeltaRow{Mengen}{G\dsep e\dsep x\dsep y}
  \DeltaRow{Funktionen}{\circ}
}

\FormulaDefDeltaK[Begriff der abelschen Gruppe]{%
  \AbelianGroupAlg{G}{\circ}{e}%
}{AbelianGroupDef}{%
  \DeltaRow{Mengen}{G\dsep e\dsep x\dsep y}
  \DeltaRow{Funktionen}{\circ}
  \DeltaRow{\textbf{Axiome}}{}
  \DeltaRow{Gruppe}{\GroupAlg{G}{\circ}{e}}
    [\FormulaRefAuto{GroupDef}[def]]
  \DeltaRow{Kommutativität}{%
    x,y\in G\vdash x\circ y=y\circ x}
    [\FormulaRefAuto{AbelianGroupCommutativityLaw}[ax]]
  \DeltaRow{\textbf{Neues Symbol}}{}
  \DeltaPrem{Abelsche Gruppe}{\AbelianGroupAlg{G}{\circ}{e}}
}

\FormulaAxiomDeltaK[Gruppenanteil einer abelschen Gruppe]{%
  \AbelianGroupAlg{G}{\circ}{e}\vdash\GroupAlg{G}{\circ}{e}%
}{AbelianGroupBaseGroupAxiom}{%
  \DeltaRow{Mengen}{G\dsep e}
  \DeltaRow{Funktionen}{\circ}
}

\FormulaAxiomDeltaK[Kommutativität in einer abelschen Gruppe]{%
  \AbelianGroupAlg{G}{\circ}{e}\dsep x,y\in G
  \vdash x\circ y=y\circ x%
}{AbelianGroupCommutativityAxiom}{%
  \DeltaRow{Mengen}{G\dsep e\dsep x\dsep y}
  \DeltaRow{Funktionen}{\circ}
}

\FormulaThmDeltaK[Abelsche Gruppen sind kommutative Monoide]{%
  \AbelianGroupAlg{G}{\circ}{e}
  \vdash\CommutativeMonoidAlg{G}{\circ}{e}%
}{AbelianGroupIsCommutativeMonoid}{%
  \DeltaRow{Mengen}{G\dsep e\dsep x\dsep y}
  \DeltaRow{Funktionen}{\circ}
}
\begin{tabproof}
  \proofstep{1}{\AbelianGroupAlg{G}{\circ}{e}}{\rA}
  \proofstep{1}{\GroupAlg{G}{\circ}{e}}{%
    \FormulaRefAuto{AbelianGroupBaseGroupAxiom}[ax]{1}}
  \proofstep{1}{\MonoidAlg{G}{\circ}{e}}{%
    \FormulaRefAuto{GroupBaseMonoidAxiom}[ax]{2}}
  \proofstep{4}{x\in G}{\rA}
  \proofstep{5}{y\in G}{\rA}
  \proofstep{1,4,5}{x\circ y=y\circ x}{%
    \FormulaRefAuto{AbelianGroupCommutativityAxiom}[ax]{1,4,5}}
  \proofstep{1}{\CommutativeMonoidAlg{G}{\circ}{e}}{%
    \FormulaRefAuto{CommutativeMonoidDef}[def]{3,6}}
\end{tabproof}

\chapter{Homomorphismen abelscher Gruppen}

Für abelsche Gruppen bleibt die Erhaltungsgleichung unverändert. Die
Kommutativität ist eine Eigenschaft von Quell- und Zielgruppe, keine weitere
Forderung an die Abbildung.

\FormulaDefDeltaK[Homomorphismen abelscher Gruppen]{%
  \AbelianGroupHom{f}{G,\circ,e}{H,\diamond,u}%
}{AbelianGroupHomDef}{%
  \DeltaRow{Mengen}{G\dsep H\dsep e\dsep u}
  \DeltaRow{Funktionen}{f\dsep\circ\dsep\diamond}
  \DeltaRow{\textbf{Axiome}}{}
  \DeltaPrem{Quellstruktur}{\AbelianGroupAlg{G}{\circ}{e}}
  \DeltaPrem{Zielstruktur}{\AbelianGroupAlg{H}{\diamond}{u}}
  \DeltaPrem{Homomorphismus}{\GroupHom{f}{G,\circ,e}{H,\diamond,u}}
  \DeltaRow{\textbf{Neues Symbol}}{}
  \DeltaPrem{Homomorphismus}{\AbelianGroupHom{f}{G,\circ,e}{H,\diamond,u}}
}

\FormulaAxiomDeltaK[Quellstruktur eines Homomorphismus abelscher Gruppen]{%
  \AbelianGroupHom{f}{G,\circ,e}{H,\diamond,u}
  \vdash \AbelianGroupAlg{G}{\circ}{e}%
}{AbelianGroupHomSourceAxiom}{%
  \DeltaRow{Mengen}{G\dsep H\dsep e\dsep u}
  \DeltaRow{Funktionen}{f\dsep\circ\dsep\diamond}
}

\FormulaAxiomDeltaK[Zielstruktur eines Homomorphismus abelscher Gruppen]{%
  \AbelianGroupHom{f}{G,\circ,e}{H,\diamond,u}
  \vdash \AbelianGroupAlg{H}{\diamond}{u}%
}{AbelianGroupHomTargetAxiom}{%
  \DeltaRow{Mengen}{G\dsep H\dsep e\dsep u}
  \DeltaRow{Funktionen}{f\dsep\circ\dsep\diamond}
}

\FormulaAxiomDeltaK[Gruppenhomomorphismusanteil eines Homomorphismus abelscher Gruppen]{%
  \AbelianGroupHom{f}{G,\circ,e}{H,\diamond,u}
  \vdash \GroupHom{f}{G,\circ,e}{H,\diamond,u}%
}{AbelianGroupHomBaseHomAxiom}{%
  \DeltaRow{Mengen}{G\dsep H\dsep e\dsep u}
  \DeltaRow{Funktionen}{f\dsep\circ\dsep\diamond}
}

\FormulaDefDeltaK[Isomorphismen abelscher Gruppen]{%
  \AbelianGroupIso{f}{G,\circ,e}{H,\diamond,u}%
}{AbelianGroupIsoDef}{%
  \DeltaRow{Mengen}{G\dsep H\dsep e\dsep u}
  \DeltaRow{Funktionen}{f\dsep\circ\dsep\diamond}
  \DeltaRow{\textbf{Axiome}}{}
  \DeltaPrem{Homomorphismus}{\AbelianGroupHom{f}{G,\circ,e}{H,\diamond,u}}
  \DeltaPrem{Bijektivität}{f\colon G\bij H}
  \DeltaRow{\textbf{Neues Symbol}}{}
  \DeltaPrem{Isomorphismus}{\AbelianGroupIso{f}{G,\circ,e}{H,\diamond,u}}
}

\FormulaAxiomDeltaK[Homomorphismusanteil eines Isomorphismus abelscher Gruppen]{%
  \AbelianGroupIso{f}{G,\circ,e}{H,\diamond,u}
  \vdash \AbelianGroupHom{f}{G,\circ,e}{H,\diamond,u}%
}{AbelianGroupIsoHomAxiom}{%
  \DeltaRow{Mengen}{G\dsep H\dsep e\dsep u}
  \DeltaRow{Funktionen}{f\dsep\circ\dsep\diamond}
}

\FormulaAxiomDeltaK[Bijektivität eines Isomorphismus abelscher Gruppen]{%
  \AbelianGroupIso{f}{G,\circ,e}{H,\diamond,u}
  \vdash f\colon G\bij H%
}{AbelianGroupIsoBijectiveAxiom}{%
  \DeltaRow{Mengen}{G\dsep H\dsep e\dsep u}
  \DeltaRow{Funktionen}{f\dsep\circ\dsep\diamond}
}

Ein Gruppenisomorphismus transportiert die Kommutativität in beiden
Richtungen. Deshalb stimmt diese Spezialisierung mit der üblichen
Isomorphie abelscher Gruppen überein.

% Die kanonische Hauptaussage bleibt in B41; Konstruktion und Tabellen in 41E.
% Die fünf Abschnittsnummern und die anonymen Formelanker bleiben erhalten.
\chapter{Formale Differenzen und der Einbettungssatz}

\section{Differenzenpaare}

Ein kommutatives Monoid besitzt nicht notwendig inverse Elemente.
Ist seine Operation kürzbar, lässt es sich jedoch durch formale
Differenzen zu einer abelschen Gruppe erweitern, ohne verschiedene
Elemente des Monoids gleichzusetzen. Die Konstruktion verallgemeinert
den Übergang von den natürlichen zu den ganzen Zahlen aus Band~17.

Sei im Folgenden \((M,\star,e)\) ein kommutatives kürzbares Monoid.
Ein Paar \((a,b)\in M\times M\) steht für die noch zu bildende
Differenz von \(a\) und \(b\). Wir vergleichen zwei solche Paare durch
\[
  (a,b)\sim(c,d)\quad\Longleftrightarrow\quad a\star d=b\star c.
\]
Diese Relation ist eine Äquivalenzrelation. Beim Nachweis der
Transitivität wird die Kürzbarkeit gebraucht. In der ausführlichen
Notation heißen der Paarträger
\(\FormalDifferenceCarrier{M,\star,e}\) und die Relation
\(\FormalDifferenceRel{M,\star,e}\).
Die \DifferencesResultLink{GroupCompletionPairEquivalence}%
  {formale Konstruktion der Äquivalenzrelation} steht in den
\DifferencesProofsLink{}.

% Neun ausgelagerte Deklarationen: section*.14 bis section*.22.
\phantomsection\phantomsection\phantomsection
\phantomsection\phantomsection\phantomsection
\phantomsection\phantomsection\phantomsection

\section{Quotient und Differenzenklassen}

Wir schreiben \([a,b]\) für die Äquivalenzklasse des Paares \((a,b)\)
und setzen
\[
  K(M):=(M\times M)/{\sim}.
\]
Dabei ist die gesamte Struktur \((M,\star,e)\) festgehalten;
\(K(M)\) ist hier die kurze Schreibweise für
\(\GroupCompletionCarrier{M,\star,e}\), und \([a,b]\) steht für
\(\FormalDifferenceClass{M,\star,e}{a}{b}\).
Jedes Element von \(K(M)\) besitzt einen solchen Repräsentanten.
Entscheidend ist das Klassenkriterium
\[
  [a,b]=[c,d]\quad\Longleftrightarrow\quad a\star d=b\star c.
\]
Die \DifferencesResultLink{GroupCompletionClassEquality}%
  {Gleichheit von Differenzenklassen} wird damit vollständig auf eine
Gleichung im Ausgangsmonoid zurückgeführt.

% Fünf ausgelagerte Deklarationen: section*.23 bis section*.27.
\phantomsection\phantomsection\phantomsection
\phantomsection\phantomsection

\section{Operation auf formalen Differenzen}

Die Operation auf dem Quotienten wird durch
\[
  [a,b]\boxplus[c,d]:=[a\star c,b\star d]
\]
festgelegt. Gleiche Eingabeklassen liefern dieselbe Ausgabeklasse;
die Vorschrift hängt also nicht von der Wahl der Repräsentanten ab.
Die \DifferencesResultLink{GroupCompletionOperationCompatible}%
  {Repräsentantenunabhängigkeit} und der
\DifferencesResultLink{GroupCompletionOperationQuotientDescent}%
  {Quotientenabstieg} werden in den Beweistabellen nachgewiesen.

Assoziativität und Kommutativität folgen aus den entsprechenden
Gesetzen in \(M\). Das neutrale Element ist \([e,e]\).
Im Hauptsatz bezeichnen
\(\GroupCompletionOp{M,\star,e}\) diese Operation und
\(\GroupCompletionZero{M,\star,e}\) das neutrale Element.

% Zehn ausgelagerte Deklarationen: section*.28 bis section*.37.
\phantomsection\phantomsection\phantomsection
\phantomsection\phantomsection\phantomsection
\phantomsection\phantomsection\phantomsection\phantomsection

\section{Kanonische Einbettung}

Die Abbildung
\[
  \iota\colon M\to K(M),\qquad \iota(a):=[a,e],
\]
erhält die Operation und das neutrale Element:
\[
  \iota(a\star b)=\iota(a)\boxplus\iota(b),\qquad
  \iota(e)=[e,e].
\]
Sie ist außerdem injektiv. Denn \([a,e]=[b,e]\) bedeutet nach dem
Klassenkriterium \(a\star e=e\star b\), also \(a=b\).
Die vollständige Bezeichnung dieser kanonischen Einbettung lautet
\(\GroupCompletionEmbed{M,\star,e}\).
Ihre \DifferencesResultLink{GroupCompletionEmbeddingMonoidHom}%
  {Homomorphie} und
\DifferencesResultLink{GroupCompletionEmbeddingInjective}%
  {Injektivität} werden in den Beweistabellen hergeleitet.

% Sechs ausgelagerte Deklarationen: section*.38 bis section*.43.
\phantomsection\phantomsection\phantomsection
\phantomsection\phantomsection\phantomsection

\section{Inverse Klassen und der Einbettungssatz}

Das Vertauschen der beiden Einträge liefert ein Inverses:
\[
  [a,b]\boxplus[b,a]=[e,e]=[b,a]\boxplus[a,b].
\]
Damit ist \(K(M)\) eine abelsche Gruppe. Zusammen mit der
injektiven strukturtreuen Abbildung \(\iota\) ergibt sich der
folgende Einbettungssatz.

% Drei ausgelagerte Hilfssätze: section*.44 bis section*.46.
% Der Hauptsatz behält 41.4.5.4 und section*.47; danach folgt section*.48.
\phantomsection\phantomsection\phantomsection
\setcounter{formulaThm}{3}
\Needspace{20\baselineskip}
\hypertarget{differences.main}{}
% Unveränderte kanonische Aussage aus Band 41; nur dort deklarieren.
\FormulaThmDeltaK[Einbettungssatz für kommutative kürzbare Monoide]{%
  \begin{aligned}[t]
    &\CommutativeMonoidAlg{M}{\star}{e}\dsep
      \mathsf{K\ddot{u}rz}(M,\star)\vdash{}\\[-2pt]
    &\AbelianGroupAlg{\GroupCompletionCarrier{M,\star,e}}
      {\GroupCompletionOp{M,\star,e}}{\GroupCompletionZero{M,\star,e}}
      \land{}\\[-2pt]
    &\CommutativeMonoidHom{\GroupCompletionEmbed{M,\star,e}}{M,\star,e}
      {\GroupCompletionCarrier{M,\star,e},
       \GroupCompletionOp{M,\star,e},\GroupCompletionZero{M,\star,e}}
      \land{}\\[-2pt]
    &\GroupCompletionEmbed{M,\star,e}\colon
      M\inj\GroupCompletionCarrier{M,\star,e}
  \end{aligned}
}{GroupCompletionEmbeddingTheorem}{%
  \DeltaRow{Mengen}{M\dsep e}
  \DeltaRow{Funktionen}{\star\dsep\GroupCompletionOp{M,\star,e}
    \dsep\GroupCompletionEmbed{M,\star,e}}
}


\noindent\emph{Beweis:} Die \DifferencesReadingLink{} erklärt die
Konstruktion und den vollständigen Beweis. Die
\DifferencesProofsLink{} enthalten alle lokalen Definitionen,
Hilfssätze und formalen Ableitungen einschließlich des
\href{\DifferencesEditionDirectory_B41-differences-proofs.pdf\#differences.proof}%
  {Schlussbeweises zum Einbettungssatz}.

Für \((M,\star,e)=(\mathbb N,+,0)\) entstehen dieselbe
Kreuzsummenrelation und dieselbe Klassenaddition wie bei den ganzen
Zahlen in Band~17. Die Kürzbarkeit begründet hier die Äquivalenzrelation;
die anschließende Gleichheitsreflexion der Einbettung folgt aus den
Neutralitätsgesetzen. Der Satz behauptet eine Einbettung und setzt
keine bereits vorhandene Subtraktion in \(M\) voraus.


\chapter{Die additive Gruppe der ganzen Zahlen}

\FormulaThmDeltaK[Additive abelsche Gruppe der ganzen Zahlen]{%
  \AbelianGroupAlg{\mathbb Z}{+}{\IntZero}%
}{IntegerAdditiveAbelianGroup}{%
  \DeltaRow{Mengen}{\mathbb Z\dsep \IntZero\dsep z\dsep w\dsep u}
  \DeltaRow{Funktionen}{+}
}
\begin{tabproofsplitwide}
  \proofpartwideindR[Additive Halbgruppe]{%
    \SemigroupAlg{\mathbb Z}{+}%
  }{\SemigroupAlg{\mathbb Z}{+}}[IntegerAdditiveSemigroup]
    \proofstepwidestar[1]{z\in\mathbb Z}{\rA}
    \proofstepwidestar[2]{w\in\mathbb Z}{\rA}
    \proofstepwidestar[1,2]{z+w\in\mathbb Z}{%
      \FormulaRefAuto{IntegerAddClosureAxiom}[ax]{1,2}}
    \proofstepwidestar[4]{u\in\mathbb Z}{\rA}
    \proofstepwidestar[1,2,4]{(z+w)+u=z+(w+u)}{%
      \FormulaRefAuto{IntegerAddAssociativeAxiom}[ax]{1,2,4}}
    \proofstepwidestar[]{\SemigroupAlg{\mathbb Z}{+}}{%
      \FormulaRefAuto{SemigroupDef}[def]{3,5}}
  \closeproofpartwideindR

  \proofpartwideindR[Additives Monoid]{%
    \MonoidAlg{\mathbb Z}{+}{\IntZero}%
  }{\MonoidAlg{\mathbb Z}{+}{\IntZero}}[IntegerAdditiveMonoid]
    \proofstepwidestar[]{\SemigroupAlg{\mathbb Z}{+}}{%
      \FormulaRefAuto{IntegerAdditiveSemigroup}[thm]}
    \proofstepwidestar[]{\IntZero\in\mathbb Z}{%
      \FormulaRefAuto{IntZeroCarrier}[thm]}
    \proofstepwidestar[3]{z\in\mathbb Z}{\rA}
    \proofstepwidestar[3]{z+\IntZero=z\land\IntZero+z=z}{%
      \FormulaRefAuto{IntegerAddZeroAxiom}[ax]{3}}
    \proofstepwidestar[3]{\IntZero+z=z}{%
      \rAEb{4}}
    \proofstepwidestar[3]{z+\IntZero=z\land\IntZero+z=z}{%
      \FormulaRefAuto{IntegerAddZeroAxiom}[ax]{3}}
    \proofstepwidestar[3]{z+\IntZero=z}{%
      \rAEa{6}}
    \proofstepwidestar[]{\MonoidAlg{\mathbb Z}{+}{\IntZero}}{%
      \FormulaRefAuto{MonoidDef}[def]{1,2,5,7}}
  \closeproofpartwideindR

  \proofpartwideindR[Additive Gruppe]{%
    \GroupAlg{\mathbb Z}{+}{\IntZero}%
  }{\GroupAlg{\mathbb Z}{+}{\IntZero}}[IntegerAdditiveGroup]
    \proofstepwidestar[]{\MonoidAlg{\mathbb Z}{+}{\IntZero}}{%
      \FormulaRefAuto{IntegerAdditiveMonoid}[thm]}
    \proofstepwidestar[2]{z\in\mathbb Z}{\rA}
    \proofstepwidestar[2]{-z\in\mathbb Z}{%
      \FormulaRefAuto{IntNegClosure}[thm]{2}}
    \proofstepwidestar[2]{%
      z+(-z)=\IntZero\land(-z)+z=\IntZero}{%
      \FormulaRefAuto{IntegerAddInverseAxiom}[ax]{2}}
    \proofstepwidestar[2]{%
      \exists w\in\mathbb Z\,
      \bigl(z+w=\IntZero\land w+z=\IntZero\bigr)}{%
      \rEI{\rAI{3,4}}}
    \proofstepwidestar[]{\GroupAlg{\mathbb Z}{+}{\IntZero}}{%
      \FormulaRefAuto{GroupDef}[def]{1,\rUI{\rRI{2,5}}}}
  \closeproofpartwideindR

  \proofpartwide{Kommutativität}
    \proofstepwidestar[]{\GroupAlg{\mathbb Z}{+}{\IntZero}}{%
      \FormulaRefAuto{IntegerAdditiveGroup}[thm]}
    \proofstepwidestar[2]{z\in\mathbb Z}{\rA}
    \proofstepwidestar[3]{w\in\mathbb Z}{\rA}
    \proofstepwidestar[2,3]{z+w=w+z}{%
      \FormulaRefAuto{IntegerAddCommutativeAxiom}[ax]{2,3}}
    \proofstepwidestar[]{%
      \AbelianGroupAlg{\mathbb Z}{+}{\IntZero}}{%
      \FormulaRefAuto{AbelianGroupDef}[def]{1,4}}
  \closeproofpartwide
\end{tabproofsplitwide}

\end{document}
